\chapter{Comunicação de Dados via Wi-Fi}
\label{cap:wifi}

Um sensor que apenas acende um LED ou imprime valores no monitor serial, como nos capítulos anteriores, já é útil para aprender — mas o que realmente distingue um projeto de \textbf{IoT} (\emph{Internet of Things}) é conectar esses dados ao resto do mundo: um painel de controle, um banco de dados, um aplicativo de celular. O \ESP{} tem um rádio Wi-Fi embutido exatamente para isso, e o framework Arduino expõe seu uso através da biblioteca \code{WiFi.h}.

\section{Conectando-se a uma rede}

\begin{codigoc}{src/main.cpp — conectando ao Wi-Fi}
#include <WiFi.h>

#define WIFI_SSID     "MinhaRede"
#define WIFI_PASSWORD "minha_senha"

void setup() {
    Serial.begin(115200);

    WiFi.begin(WIFI_SSID, WIFI_PASSWORD);
    Serial.print("Conectando ao Wi-Fi");

    while (WiFi.status() != WL_CONNECTED) {
        delay(500);
        Serial.print(".");
    }

    Serial.println();
    Serial.print("Conectado! Endereco IP: ");
    Serial.println(WiFi.localIP());
}

void loop() {
    // ...
}
\end{codigoc}

\code{WiFi.begin(ssid, senha)} inicia a conexão de forma assíncrona; \code{WiFi.status()} permite consultar seu andamento a qualquer momento, retornando \code{WL\_CONNECTED} assim que a conexão é estabelecida. \code{WiFi.localIP()} retorna o endereço IP obtido na rede local — útil para depuração e, como veremos no \cref{cap:ota}, para atualizações remotas.

\begin{nota}
O laço de espera com \code{delay(500)} dentro de \code{setup()} é uma das poucas exceções aceitáveis à regra de ``nunca bloquear'' apresentada no \cref{cap:interrupcoes}: como acontece uma única vez, durante a inicialização, antes de \code{loop()} sequer começar a rodar, bloquear a execução aqui não compromete a responsividade do sistema em operação normal.
\end{nota}

\section{Enviando dados: requisições HTTP}

Com a rede conectada, a forma mais simples de enviar dados para um servidor é uma requisição \textbf{HTTP}, usando a biblioteca \code{HTTPClient.h}:

\begin{codigoc}{Enviando uma leitura de sensor via HTTP POST}
#include <WiFi.h>
#include <HTTPClient.h>

#define URL_SERVIDOR "http://exemplo.com/api/leituras"

void enviar_leitura(float temperatura) {
    if (WiFi.status() != WL_CONNECTED) {
        return;   // sem rede, nao ha o que enviar
    }

    HTTPClient http;
    http.begin(URL_SERVIDOR);
    http.addHeader("Content-Type", "application/json");

    char corpo[64];
    snprintf(corpo, sizeof(corpo), "{\"temperatura\": %.1f}", temperatura);

    int codigo_resposta = http.POST(corpo);
    Serial.print("[HTTP] Resposta do servidor: ");
    Serial.println(codigo_resposta);

    http.end();   // libera os recursos da requisicao
}
\end{codigoc}

\begin{dica}
Repare no uso de \code{snprintf} para montar o corpo da requisição em um \code{char} de tamanho fixo, em vez de concatenar strings dinamicamente com o tipo \code{String} do Arduino. Concatenação repetida de \code{String} realoca memória no heap a cada operação — inofensivo em um exemplo isolado, mas uma fonte conhecida de \textbf{fragmentação de memória} em firmwares que rodam por dias ou semanas sem reiniciar. \code{snprintf} — a mesma função de formatação apresentada no \cref{cap:entrada-saida}, aqui aplicada a um buffer em vez da tela — evita esse problema por completo.
\end{dica}

O valor retornado por \code{http.POST(...)} é o código de status HTTP da resposta (\code{200} para sucesso, \code{404} se o endereço não existe, \code{500} para erro no servidor, e assim por diante) — ou um valor negativo se a própria requisição falhou (por exemplo, servidor inacessível).

\section{Lidando com quedas de conexão}

Redes Wi-Fi não são perfeitamente confiáveis: a conexão pode cair a qualquer momento. Verificar \code{WiFi.status() != WL\_CONNECTED} antes de cada envio, como no exemplo acima, evita que o firmware trave tentando enviar dados por uma rede indisponível. Para tentar reconectar automaticamente, sem bloquear o restante do laço principal, a mesma técnica de temporização por \code{millis()} do \cref{cap:interrupcoes} se aplica:

\begin{codigoc}{Reconexão não bloqueante}
unsigned long ultima_tentativa = 0;
const unsigned long INTERVALO_RECONEXAO = 10000;   // 10 segundos

void verificar_wifi() {
    if (WiFi.status() == WL_CONNECTED) {
        return;
    }

    unsigned long agora = millis();
    if (agora - ultima_tentativa >= INTERVALO_RECONEXAO) {
        ultima_tentativa = agora;
        Serial.println("Wi-Fi desconectado. Tentando reconectar...");
        WiFi.begin(WIFI_SSID, WIFI_PASSWORD);
    }
}
\end{codigoc}

Chamar \code{verificar\_wifi()} a cada volta de \code{loop()} garante que o firmware tente se reconectar periodicamente, sem nunca travar o restante do programa esperando por uma rede que pode não voltar tão cedo.

\section{Além do HTTP: uma palavra sobre MQTT}

Requisições HTTP funcionam bem para envios pontuais e periódicos, como no exemplo desta seção. Para projetos de IoT com muitos dispositivos trocando mensagens constantemente — sensores publicando leituras, atuadores recebendo comandos, tudo em tempo quase real — o protocolo mais usado da indústria é o \textbf{MQTT}: um protocolo de mensagens leve, baseado em um modelo de publicação/assinatura (\emph{publish/subscribe}) através de um servidor central chamado \emph{broker}. Ele não faz parte desta apostila introdutória, mas é o passo natural seguinte para quem for construir um projeto de IoT com múltiplos dispositivos — bibliotecas como a \code{PubSubClient}, compatíveis com o framework Arduino, tornam sua adoção direta.

\begin{esp32box}
O rádio Wi-Fi é, disparadamente, o componente que mais consome energia no \ESP{} — voltando ao \cref{cap:energia}, um dispositivo que precisa economizar bateria e também reportar dados pela rede normalmente segue o padrão: acordar do \emph{deep sleep}, conectar-se ao Wi-Fi, enviar os dados pendentes, desconectar, e voltar a dormir — minimizando o tempo em que o rádio permanece ativo, em vez de mantê-lo conectado continuamente.
\end{esp32box}

Com o \ESP{} conectado à rede, o próximo capítulo mostra como usar essa mesma conectividade para atualizar o firmware remotamente — sem precisar de um cabo USB.
